% Source: https://artofproblemsolving.com/community/c3h1754998_probability_problem
% Stable identity: aops-topic-1754998-q01
The point $(x_1,\dots,x_n)$ is uniform in the unit cube $[0,1]^n$, of volume $1$, so the probability is the volume of
\[
T=\{x_i\ge0,\ x_1+\cdots+x_n<1\}.
\]
Consider the partial sums $s_k=x_1+\cdots+x_k$. The map $(x_1,\dots,x_n)\mapsto(s_1,\dots,s_n)$ is linear with a triangular matrix having $1$s on the diagonal, so it preserves volume, and it maps $T$ bijectively onto
\[
T'=\{0\le s_1\le s_2\le\cdots\le s_n<1\}
\]
(the inverse is $x_1=s_1$, $x_k=s_k-s_{k-1}$).

Now the cube $[0,1)^n$ is the union of the $n!$ regions $\{s_{\sigma(1)}\le s_{\sigma(2)}\le\cdots\le s_{\sigma(n)}\}$, one for each permutation $\sigma$ of the coordinates; two of them overlap only where two coordinates are equal, a set of volume $0$. Permuting coordinates preserves volume, so all $n!$ regions have the same volume, which is therefore $\dfrac{1}{n!}$. One of them is $T'$, so
\[
P(x_1+\cdots+x_n<1)=\operatorname{vol}(T)=\operatorname{vol}(T')=\boxed{\dfrac{1}{n!}} .\qquad\blacksquare
\]
(For $n=2$ this is the triangle below the line $x_1+x_2=1$, of area $\dfrac{1}{2}$; for $n=3$ the corner tetrahedron, of volume $\dfrac{1}{6}$. The other suggestion, $\dfrac{1}{n(n-1)}$, agrees with $\dfrac{1}{n!}$ only for $n=2,3$; for $n=4$ it gives $\dfrac{1}{12}\ne\dfrac{1}{24}$.)
